\section{Discussion}
\label{sec:discussion}

\subsection{Key Findings}

\textbf{The domain content differentiation result is the primary empirical contribution.} $\eta^2 > 0.91$ means The Pile's domains are categorically, not marginally, distinct in cognitive content. Prior data mixing research (RegMix, DoReMi, DCLM) has assumed this differentiation without measuring it. We measure it.

\textbf{The trajectory non-uniformity result enables the next study.} The infrastructure validation --- \texttt{doc\_idx.npy} identity mapping, \texttt{step\_to\_sample()}, JSONL.zst streaming (3.5 min for 600K docs) --- constitutes the missing link between Pythia's training procedure and domain exposure analysis.

\textbf{The 70M reversal requires replication, not rejection.} Three competing explanations, in decreasing plausibility: (1) 70M MMLU floor effect --- MMLU is near-random at 70M, making $\rho(\text{anything, MMLU})$ noise-dominated [HIGH plausibility]; (2) $N=10$ checkpoint sampling artifact [HIGH plausibility]; (3) Pile-CC multicollinearity confounding Wikipedia's univariate Spearman [MEDIUM plausibility]; (4) genuine domain-benchmark reversal at all scales [LOW plausibility]. The 1B and 6.9B results will adjudicate.

\subsection{Limitations}

\textbf{L1:} Six Pile domains (Books3, OWT2, GitHub, OpenSubtitles, BookCorpus2, YoutubeSubtitles) have zero exposure in the first-shard sample. Fix: full 134M-document lookup across all 30 shards ($\sim$150 compute-hours).

\textbf{L2:} Evaluation cache incomplete at analysis time (70M: 10/154, 1B: 2/154, 6.9B: 0/154). Fix: complete the running evaluation pipeline (462 total evaluations).

\textbf{L3:} h-m1 used domain-representative generated text rather than real Pile documents; $\eta^2$ may be inflated. Directional claims are robust; quantitative values are upper bounds.

\textbf{L4:} Observational design --- correlations do not establish causal domain effects on benchmark performance.

\subsection{Broader Impact}

This work contributes infrastructure for pre-training data attribution research and documents a reproducibility concern for future work using partial Pile samples. No negative societal impacts are anticipated from this analytical methodology.
